\Needspace{22\baselineskip}\vspace{6mm}
\pdfbookmark[1]{Practical Instructions for Making Artificial Magnets.}{fischer-magnete}
\begin{center}\small Göttingische gelehrte Anzeigen.\quad 1832.\quad September 10.\end{center}

\textit{Practical instructions for making and assembling artificial magnets to best advantage, especially horseshoes, straight bars, compass needles and other needles, together with the latest discovery for giving them the greatest attractive force, for naturalists, physicians, seafarers, technicians and all other kinds of metalworkers, by \textup{\textsc{Friedrich Fischer}}, teacher and practising technician. Heilbronn.\ \textup{1833.} At Class’s bookshop.} (58 pages in 8vo. With two lithographed plates.)

Of the many phenomena exhibited by magnetic force, artificial magnets carrying heavy loads especially excite the admira\SourceLeaf{598}tion of enthusiasts, the magnets weighing only a very small fraction of their loads. It is of only secondary scientific interest, however, and barely rates a few words in the most comprehensive recent work on physics. Designing the artificial magnets in horseshoe form best suited to that purpose is therefore seen less as the physicist’s business than as a matter for people who make a living from it and sometimes offer for sale supposedly new methods of their own under a seal of secrecy. A recent case of this kind prompted the author of the present short work to make experiments of his own, and it is to his credit that he publishes the results without reserve or secretiveness. By far the greater part of the work is devoted to artificial horseshoe magnets. Their most advantageous proportions, the selection and treatment of the steel, and the stroking methods to be used are described in an intelligible way. Enthusiasts can be assured that by following the instructions they can always obtain such magnets with a very high load-bearing capacity.

The reviewer would confine himself to recommending the present work were it not that the very stroking methods used by the author prompt a few remarks of his own. These methods are presumably what is meant by the latest discovery mentioned in the title, although the author himself admits that they may already be known elsewhere. The remarks may be regarded as a supplement of some importance to \textsc{Coulomb}’s experiments on the methods of \textsc{Knight}, \textsc{Duhamel}, \textsc{Michell}, \textsc{Canton} and \textsc{Aepinus}, all well known to physicists.

To induce magnetism in a horseshoe magnet that is being made, the author uses an existing magnet of the same shape. His procedure consists of two operations performed in succession, in both of which the two poles always stroke at the same time. In the first, one pole of the stroking magnet follows the other along its path. In the second, one pole is guided along one arm and the other along the other, from the bend towards the end. There is no need to specify the requisite arrangement of the poles here. Before the second operation the author also advises warming the magnet to be stroked and continuing the work until it has cooled.

\SourceLeaf{599}It is easy to see that the first operation is exactly the double touch invented by \textsc{Michell}. The second operation, by contrast, agrees in essentials with \textsc{Duhamel}’s procedure, except that the separate straight bars (or bundles of bars) used by \textsc{Duhamel} for stroking offer some advantages for producing stronger magnetism that are lost when a single horseshoe magnet is used (especially because one cannot start at the middle of the bend). As is well known, \textsc{Canton}’s method consists solely of combining the methods of \textsc{Michell} and \textsc{Duhamel}. The author’s procedure is thus essentially \textsc{Canton}’s with the modifications that naturally come with using a horseshoe-shaped stroking magnet. Its only genuinely new feature is the preliminary warming, whose effectiveness will probably need further confirmation. In some experiments of his own, the reviewer found no particular effect from it at all.

What especially deserves notice here, however, is that physicists, following \textsc{Coulomb}’s lead, do not consider \textsc{Canton}’s method any improvement at all. In that celebrated physicist’s judgment the intensity of the magnetism induced is always determined solely by the method applied last, so preceding it with \textsc{Michell}’s way of stroking is entirely superfluous. On the other hand, it is clear from our author’s remarks that he regards the combination of his two operations as essential. The reviewer readily acknowledges that these remarks, which bear the stamp of an unassuming devotion to truth, first led him to doubt the general validity of the principle underlying \textsc{Coulomb}’s judgment. A great many experiments using distinctive means of testing, which afford the greatest precision and will be described elsewhere, have fully justified this doubt.

As is well known, the improvement to \textsc{Michell}’s stroking method that we owe to \textsc{Aepinus} is by far the most effective. With somewhat thicker pieces of steel, every other method, including \textsc{Duhamel}’s, falls well short of it. \textsc{Coulomb}’s experiments have put this beyond all doubt. Physicists therefore use \textsc{Aepinus}’s method exclusively to magnetise such pieces of steel as strongly as possible. The result consistently obtained in the reviewer’s experiments mentioned above is therefore striking, and unexpected on the assump\SourceLeaf{600}tions accepted until now. Steel magnetised as strongly as possible by \textsc{Aepinus}’s method invariably gains considerable strength when it is \textit{afterwards} stroked repeatedly by \textsc{Duhamel}’s procedure, even though the latter alone can produce considerably less strength than \textsc{Aepinus}’s method. The reviewer confines himself here to reporting the fact and leaves aside the attempt at an explanation, which in any case suggests itself fairly readily. Although these observations were made directly only in magnetising straight bars, there can be no doubt that the combination of \textsc{Aepinus}’s and \textsc{Duhamel}’s methods must likewise induce the strongest possible magnetism in horseshoe-shaped laminations. Some special arrangements are then needed if the combination is to be carried out conveniently in its unmodified form. Anyone who is unwilling to make those arrangements or has no suitable straight bars at hand will still obtain very great strength, even with somewhat thicker laminations, by using a horseshoe-shaped stroking magnet as the author directs. The strength will fall short of the maximum attainable. The magnet is also quite easy to handle.

What the author says about magnetising straight bars is limited to the operations to be performed if one wants to stroke them with a horseshoe magnet. This gives considerable strength, though the sequence of \textsc{Aepinus}’s and \textsc{Duhamel}’s methods mentioned above gives still greater strength and leaves nothing to be desired in ease of use either.

The author recommends storing magnetised straight bars by laying them on top of one another with like poles together. This is entirely inadvisable if they are to retain as much of their strength as possible. It is best to lay them side by side in pairs a short distance apart, with unlike poles adjacent, and to place keepers of very soft iron of suitable length against them.

